% Generated draft: requires mathematical and source review.
\documentclass[11pt]{article}
\usepackage[margin=25mm]{geometry}
\usepackage{fontspec}
\setmainfont{Latin Modern Roman}
\usepackage{amsmath,amssymb,mathtools,unicode-math}
\setmathfont{Latin Modern Math}
\usepackage{enumitem,needspace,xurl,hyperref}
\setlength{\parindent}{0pt}
\setlength{\parskip}{5pt}
\setlength{\emergencystretch}{3em}
\newcommand{\N}{\mathbb{N}}
\newcommand{\Z}{\mathbb{Z}}
\newcommand{\Q}{\mathbb{Q}}
\newcommand{\R}{\mathbb{R}}
\newcommand{\C}{\mathbb{C}}
\newcommand{\F}{\mathbb{F}}
\newcommand{\A}{\mathbb{A}}
\newcommand{\PP}{\mathbb{P}}
\newcommand{\E}{\mathbb{E}}
\newcommand{\eps}{\varepsilon}
\newcommand{\dd}{\,\mathrm{d}}
\begin{document}
\section*{UnsolvedMath ID: 1458}
\textbf{Source problem code:} catalogue entry
\subsection{Definitions and mathematical statement}
An abstract group \(G\) is simple if it is nontrivial and its only normal subgroups are \(\{1\}\) and \(G\). Let \(\operatorname{Th}(G)\) denote its complete first-order theory in the pure group language. A theory is stable if it has no formula with the order property: there is no formula \(\varphi(x;y)\) and tuples \((a_i)_{i\in\Z_{\geq0}}\), \((b_j)_{j\in\Z_{\geq0}}\) in a model for which
\[
\varphi(a_i;b_j)\quad\Longleftrightarrow\quad i<j.
\]
An algebraic group over a field \(K\) is an algebraic variety over \(K\) with regular multiplication and inversion; its \(K\)-rational points form an abstract group, denoted \(\mathbf H(K)\).

The statement asks whether every simple group with stable first-order theory is an algebraic group over an algebraically closed field. The background specifically discusses infinite simple groups. Under that background restriction, the pure-group-language interpretation, and an abstract-isomorphism reading of ``is an algebraic group,'' the candidate assertion is
\[
\forall G\text{ infinite simple},\quad
\operatorname{Th}(G)\text{ stable}
\Longrightarrow
\exists K\text{ algebraically closed}\ \exists\mathbf H/K\text{ an algebraic group},\quad
G\cong\mathbf H(K).
\]
The source does not say whether this is merely an abstract isomorphism or a definable identification in an ambient structure, nor whether \(\mathbf H\) must be connected, affine, or simple as an algebraic group. Those refinements are not silently imposed. The short statement also lacks the background's infinitude restriction, and therefore includes finite simple groups unless that restriction is adopted explicitly.

The supplied triage dated 2026-08-17 stresses that the stable-theory wording differs from the standard Cherlin--Zilber finite-Morley-rank algebraicity conjecture. Finite Morley rank is an additional model-theoretic rank restriction, not a synonym for stability. The rank-restricted conjecture is not substituted for the broader source wording. The triage labels the row partially solved, records important even/mixed-type results in the finite-Morley-rank setting, and reports no verified theorem covering all simple stable groups. The registry status is open. These are supplied dated assessments, not an independently verified status of every possible formulation.
\subsection{Short English statement}
The source asks whether a simple group with stable first-order theory must arise as the points of an algebraic group over an algebraically closed field. It must clarify infinitude, the language and ambient structure, the meaning of algebraicity, and whether the intended hypothesis is stability or the stronger finite-Morley-rank condition.
\subsection{Sources}
\begin{itemize}
\item Local UnsolvedMath registry, permanent ID 1458, source code catalogue entry, ``Cherlin-Zilber Conjecture,'' including supplied literature triage dated 2026-08-17 and attribution through the conjecture name to Cherlin and Zilber.
\item K. Tent, ``From the Cherlin-Zilber Conjecture via sharply 2-transitive groups to the Burnside problem,'' arXiv:2606.18207 (2026): \url{https://arxiv.org/abs/2606.18207}. The supplied triage cites this as a review of the unresolved algebraicity programme.
\end{itemize}
\end{document}
